\section*{NeurIPS Paper Checklist}

\begin{enumerate}

\item {\bf Claims}
    \item[] Question: Do the main claims made in the abstract and introduction accurately reflect the paper's contributions and scope?
    \item[] Answer: \answerYes{}
    \item[] Justification: The abstract and Section~\ref{sec:intro} list four contributions
    (frozen-ESM2 lift, edge-type ablation, GNNExplainer motif validation,
    reproducible Make-driven pipeline), each substantiated by quantitative
    results in Section~\ref{sec:results}.

\item {\bf Limitations}
    \item[] Question: Does the paper discuss the limitations of the work performed by the authors?
    \item[] Answer: \answerYes{}
    \item[] Justification: Section~\ref{sec:limitations} discusses the single-seed
    evaluation, the fixed 8\,\AA{} contact threshold, the loss of
    AlphaFold-less proteins, the modest random search budget, and the failure
    of GNNExplainer to recover ER and plastid motifs.

\item {\bf Theory assumptions and proofs}
    \item[] Question: For each theoretical result, does the paper provide the full set of assumptions and a complete (and correct) proof?
    \item[] Answer: \answerNA{}
    \item[] Justification: The paper is purely empirical and presents no theorems
    or formal claims requiring proof.

\item {\bf Experimental result reproducibility}
    \item[] Question: Does the paper fully disclose all the information needed to reproduce the main experimental results of the paper to the extent that it affects the main claims and/or conclusions of the paper (regardless of whether the code and data are provided or not)?
    \item[] Answer: \answerYes{}
    \item[] Justification: All splits are taken from the public DeepLoc 2.0 release,
    every hyperparameter and the random-search space are listed in
    Section~\ref{sec:experiments}, the chosen headline configuration is given
    in Table~\ref{tab:hyper}, and the entire pipeline is exposed as
    \texttt{make gnn-*} targets with fixed seed (\texttt{--seed 0}).

\item {\bf Open access to data and code}
    \item[] Question: Does the paper provide open access to the data and code, with sufficient instructions to faithfully reproduce the main experimental results, as described in supplemental material?
    \item[] Answer: \answerYes{}
    \item[] Justification: Code is committed to the project repository under
    \texttt{halo/gnn/} and \texttt{halo/spec/}; a \texttt{Makefile} drives every
    stage end-to-end; data are the publicly released DeepLoc~2.0 CSVs and
    AlphaFold structures fetched from AFDB by UniProt accession.

\item {\bf Experimental setting/details}
    \item[] Question: Does the paper specify all the training and test details (e.g., data splits, hyperparameters, how they were chosen, type of optimizer) necessary to understand the results?
    \item[] Answer: \answerYes{}
    \item[] Justification: Section~\ref{sec:experiments} specifies the DeepLoc
    train/val/test partition, the BCE-with-logits loss with positive-class
    re-weighting, the AdamW optimizer with cosine schedule, AMP, batch size,
    epoch budget, and the random-search space over GAT depth, width, heads,
    dropout, and learning rate. Selection is by macro val ROC-AUC; the test
    set is touched once per row.

\item {\bf Experiment statistical significance}
    \item[] Question: Does the paper report error bars suitably and correctly defined or other appropriate information about the statistical significance of the experiments?
    \item[] Answer: \answerNo{}
    \item[] Justification: Each ablation row is trained with a single fixed seed.
    A multi-seed estimate would require roughly $5\times$ the GPU-hours we had
    available on a single laptop RTX~4070 (8\,GB). We instead report
    consistency across the eight random-search configurations
    (Figure~\ref{fig:sweep}), which span a $\sim\!0.7$ point val ROC-AUC range
    and act as an architectural sensitivity check.

\item {\bf Experiments compute resources}
    \item[] Question: For each experiment, does the paper provide sufficient information on the computer resources (type of compute workers, memory, time of execution) needed to reproduce the experiments?
    \item[] Answer: \answerYes{}
    \item[] Justification: All training and explanation runs were performed on a
    single NVIDIA RTX~4070 laptop GPU (8\,GB VRAM); approximate wall-clock
    budgets per stage are given in Section~\ref{sec:experiments}. Total
    project compute is on the order of two GPU-days plus two CPU-evenings of
    AlphaFold downloads and ESM2 caching.

\item {\bf Code of ethics}
    \item[] Question: Does the research conducted in the paper conform, in every respect, with the NeurIPS Code of Ethics?
    \item[] Answer: \answerYes{}
    \item[] Justification: We use only publicly released resources (DeepLoc~2.0,
    ESM2, AlphaFoldDB), perform no human-subjects research, and release no
    high-risk artefacts.

\item {\bf Broader impacts}
    \item[] Question: Does the paper discuss both potential positive societal impacts and negative societal impacts of the work performed?
    \item[] Answer: \answerYes{}
    \item[] Justification: A brief broader-impacts paragraph in
    Section~\ref{sec:discussion} notes the positive value of accurate
    localization predictors in functional annotation and drug-target
    triage, and the risk that overly confident predictions on
    out-of-distribution proteins (e.g.\ designed sequences) could mislead
    downstream wet-lab decisions.

\item {\bf Safeguards}
    \item[] Question: Does the paper describe safeguards that have been put in place for responsible release of data or models that have a high risk for misuse?
    \item[] Answer: \answerNA{}
    \item[] Justification: We release a small classifier head on top of frozen
    public embeddings; the artefact has no plausible misuse path and poses no
    high-risk dual-use concern.

\item {\bf Licenses for existing assets}
    \item[] Question: Are the creators or original owners of assets used in the paper properly credited and are the license and terms of use explicitly mentioned and properly respected?
    \item[] Answer: \answerYes{}
    \item[] Justification: DeepLoc~2.0 (academic use), ESM2 (MIT licence),
    AlphaFoldDB (CC-BY 4.0), \texttt{torch\_geometric} (MIT), and
    \texttt{scikit-learn} (BSD-3) are cited in
    Section~\ref{sec:related} and Section~\ref{sec:experiments}.

\item {\bf New assets}
    \item[] Question: Are new assets introduced in the paper well documented and is the documentation provided alongside the assets?
    \item[] Answer: \answerNA{}
    \item[] Justification: The paper does not release a new dataset or pretrained
    model. The trained GAT checkpoint and the GNNExplainer attributions are
    project artefacts retained in the repository for reproducibility but are
    not packaged as a public release.

\item {\bf Crowdsourcing and research with human subjects}
    \item[] Question: For crowdsourcing experiments and research with human subjects, does the paper include the full text of instructions given to participants and screenshots, if applicable, as well as details about compensation (if any)?
    \item[] Answer: \answerNA{}
    \item[] Justification: The work involves no crowdsourcing and no human subjects.

\item {\bf Institutional review board (IRB) approvals or equivalent for research with human subjects}
    \item[] Question: Does the paper describe potential risks incurred by study participants, whether such risks were disclosed to the subjects, and whether Institutional Review Board (IRB) approvals (or an equivalent approval/review based on the requirements of your country or institution) were obtained?
    \item[] Answer: \answerNA{}
    \item[] Justification: The work involves no human subjects and therefore did
    not require IRB review.

\item {\bf Declaration of LLM usage}
    \item[] Question: Does the paper describe the usage of LLMs if it is an important, original, or non-standard component of the core methods in this research?
    \item[] Answer: \answerYes{}
    \item[] Justification: The frozen ESM2-150M protein language model is a core
    component of the feature extractor and is described in
    Section~\ref{sec:method}. No general-purpose LLM is used in the core
    methodology; LLM-assisted writing was limited to copy-editing.

\end{enumerate}
